% ============================================================
%  Bewusst Essen — Praktischer Leitfaden zu Lebensmittelzusatzstoffen
%  faustobe.it | 2026 | CC BY-NC 4.0
% ============================================================
\documentclass[a5paper,12pt,oneside,openany]{book}

\usepackage[T1]{fontenc}
\usepackage[utf8]{inputenc}
\usepackage{ebgaramond}
\usepackage[a5paper, top=2cm, bottom=3cm, inner=2.2cm, outer=3cm]{geometry}
\usepackage[table]{xcolor}
\usepackage[ngerman]{babel}

% Colors
\definecolor{verdescuro}{HTML}{2d6a4f}
\definecolor{verdecharo}{HTML}{52b788}
\definecolor{grigioscuro}{HTML}{2b2d42}
\definecolor{gelb}{HTML}{d4a017}
\definecolor{rot}{HTML}{c1121f}
\definecolor{blauinfo}{HTML}{3a86ff}

\usepackage{tcolorbox}
\tcbuselibrary{skins,breakable}

\usepackage{booktabs}
\usepackage{tabularx}
\newcolumntype{L}{>{\raggedright\arraybackslash}X}

\usepackage{hyperref}
\hypersetup{
  colorlinks=true,
  linkcolor=verdescuro,
  urlcolor=verdescuro,
  pdftitle={Bewusst Essen},
  pdfauthor={faustobe},
}

\usepackage{qrcode}
\usepackage{fancyhdr}
\usepackage{titlesec}
\usepackage{microtype}

% Colored boxes
\tcbset{
  kastenbase/.style={
    breakable, enhanced, arc=4pt, boxrule=0.8pt,
    fonttitle=\bfseries\small,
    before upper={\parskip=4pt},
  }
}

\newcommand{\boxgruen}[2]{%
  \begin{tcolorbox}[kastenbase,
    colback=verdescuro!8!white, colframe=verdescuro, title={#1}]
  #2\end{tcolorbox}}

\newcommand{\boxgelb}[2]{%
  \begin{tcolorbox}[kastenbase,
    colback=gelb!12!white, colframe=gelb!80!black, title={#1}]
  #2\end{tcolorbox}}

\newcommand{\boxrot}[2]{%
  \begin{tcolorbox}[kastenbase,
    colback=rot!8!white, colframe=rot, title={#1}]
  #2\end{tcolorbox}}

\newcommand{\boxinfo}[2]{%
  \begin{tcolorbox}[kastenbase,
    colback=blauinfo!8!white, colframe=blauinfo, title={#1}]
  #2\end{tcolorbox}}

% Page style
\pagestyle{fancy}
\fancyhf{}
\fancyfoot[C]{\footnotesize\textcolor{grigioscuro}{\textit{Bewusst Essen\ \textbar\ faustobe.it}}}
\fancyfoot[R]{\footnotesize\thepage}
\renewcommand{\headrulewidth}{0pt}
\renewcommand{\footrulewidth}{0.4pt}
\fancypagestyle{plain}{%
  \fancyhf{}%
  \fancyfoot[C]{\footnotesize\textcolor{grigioscuro}{\textit{Bewusst Essen\ \textbar\ faustobe.it}}}%
  \fancyfoot[R]{\footnotesize\thepage}%
  \renewcommand{\headrulewidth}{0pt}%
  \renewcommand{\footrulewidth}{0.4pt}%
}

% Chapter title style
\titleformat{\chapter}[display]
  {\normalfont\LARGE\bfseries\color{verdescuro}}
  {\chaptertitlename\ \thechapter}{10pt}{\Huge}
\titlespacing*{\chapter}{0pt}{-20pt}{30pt}

% ============================================================
\begin{document}

% -------- UMSCHLAG --------
\begin{titlepage}
\pagecolor{verdescuro}
\color{white}
\vspace*{2.5cm}
\begin{center}
  {\fontsize{38}{44}\selectfont\bfseries Bewusst Essen}\\[1.2cm]
  \textcolor{verdecharo}{\rule{6cm}{0.6pt}}\\[1.2cm]
  {\large\itshape Praktischer Leitfaden zu Lebensmittelzusatzstoffen}\\[0.6cm]
  {\normalsize Mit E-Codes Reader f\"{u}r Android}\\[3cm]
  {\normalsize\textcolor{verdecharo}{faustobe.it}}
\end{center}
\vfill
\begin{center}
  {\small Version 1.0\ \textbullet\ 2026\ \textbullet\ CC BY-NC 4.0}
\end{center}
\end{titlepage}
\nopagecolor

\frontmatter
\tableofcontents
\mainmatter

% ============================================================
\chapter{So nutzen Sie die App f\"{u}r bewusstes Einkaufen}

Lernen Sie, die am st\"{a}rksten verarbeiteten Lebensmittel zu erkennen und
wirklich einfache, nat\"{u}rliche zu w\"{a}hlen --- in wenigen Minuten,
ohne Experte f\"{u}r Zusatzstoffe zu werden.

\section{Vor dem Scannen: drei schnelle \"{U}berpr\"{u}fungen}

Schauen Sie sich die Verpackung an, noch bevor Sie die App \"{o}ffnen.
Vorderseiten-Etiketten sind oft so gestaltet, dass Produkte ges\"{u}nder
wirken, als sie wirklich sind.

\begin{itemize}
  \item \textbf{Schauen Sie auf die Vorderseite der Verpackung.} Achten Sie
    auf Aussagen wie ``zuckerfrei, aber mit S\"{u}{\ss}ungsmitteln'',
    ``fettfrei, aber mit Emulgatoren'', ``Extra-Geschmack'' oder
    ``aromatisiert''. Diese Hinweise deuten oft auf ein stark ver\"{a}ndertes
    Produkt hin.
  \item \textbf{\"{O}ffnen Sie die App und scannen Sie das Produkt.} Verwenden
    Sie die Kamera, um den Barcode zu scannen, oder laden Sie ein Foto des
    Etiketts hoch. Die App zeigt Ihnen sofort die Zutatenliste und die
    enthaltenen Zusatzstoffe an.
  \item \textbf{\"{U}berpr\"{u}fen Sie drei Schl\"{u}sselpunkte.} Wie viele
    Zusatzstoffe sind aufgef\"{u}hrt (vor allem E-Nummern)? Wie lang und
    komplex ist die Zutatenliste? Wie viel Zucker, Salz und Fett wird angegeben?
\end{itemize}

Wenn Sie auf mehr als eine Frage mit ``viele'' oder ``lang/kompliziert''
antworten, haben Sie es wahrscheinlich mit einem stark verarbeiteten
Lebensmittel zu tun.

\section{Drei Regeln, um sofort zu erkennen, ob ein Lebensmittel ``zu
verarbeitet'' ist}

\textbf{Wenige Zutaten = besser.} Wenn die Liste kurz ist und vertraute Namen
enth\"{a}lt (z.\,B. ``Mehl, Wasser, Salz, \"{O}l''), ist das Produkt wenig
verarbeitet und in der Regel zuverl\"{a}ssiger.

\textbf{Viele ``unaussprechliche'' Zutaten.} Begriffe wie ``modifizierte
St\"{a}rken'', ``isolierte Proteine'', ``Mono- und Diglyceride von
Fetts\"{a}uren'', ``Geschmacksverst\"{a}rker'' oder ``komplexe Aromen''
sind ein eindeutiges Zeichen f\"{u}r ein stark verarbeitetes Lebensmittel.

\textbf{Viele E-Nummern und Zusatzstoffe.} Nicht alle E-Nummern sind
sch\"{a}dlich, aber ihr massives Vorkommen --- Farbstoffe, Konservierungsstoffe,
Stabilisatoren, Emulgatoren --- weist auf einen hohen Grad industrieller
Verarbeitung hin.

\section{So lesen Sie das Ergebnis der App}

Jedes Produkt erh\"{a}lt eine Schnellbewertung: stellen Sie sich diese Farben
wie eine Lebensmittelampel vor.

\boxgruen{Wenig verarbeitet --- gr\"{u}nes Licht}{Wenige Zutaten, wenige oder
keine E-Zusatzstoffe. Moderate Mengen an Zucker und Salz. Ausgezeichnete Wahl
f\"{u}r den t\"{a}glichen Einkauf.}

\boxgelb{M\"{a}{\ss}ig verarbeitet --- in Ma{\ss}en verwenden}{Einige E-Nummern,
nat\"{u}rliche Aromen oder Zuckerzus\"{a}tze. Gelegentlich in Ordnung, aber
nicht als Grundlage jeder Mahlzeit.}

\boxrot{Stark verarbeitet --- einschr\"{a}nken oder meiden}{Viele Zusatzstoffe,
Zuckerzus\"{a}tze, komplexe Aromen und industriell verarbeitete Fette. Besser
als seltene Ausnahme behandeln, besonders f\"{u}r Kinder.}

\section{So nutzen Sie die App praktisch beim Einkaufen}

\textbf{1.~Beginnen Sie mit dem K\"{u}hl- und Fr\"{u}hst\"{u}cksbereich.}
Scannen Sie Joghurt, Brot, Milch, Getr\"{a}nke. Suchen Sie nach Produkten mit
wenigen Zutaten, wenigen E-Nummern und einfachen Zuckern (Obst, Milch, Rohrzucker)
statt Sirupen und S\"{u}{\ss}ungsmitteln.

\textbf{2.~Dann weiter zu S\"{u}{\ss}waren und Snacks.} Wenn die Liste sehr lang
und voller E-Nummern ist, zeigt die App Rot an. W\"{a}hlen Sie dann immer die
einfachere Alternative: Obst, Vollkornbrot, N\"{u}sse.

\textbf{3.~Nutzen Sie die App zum Vergleich \"{a}hnlicher Produkte.} Oft ist das
mit weniger Zutaten nur geringf\"{u}gig unscheinbarer im Regal, aber deutlich
nat\"{u}rlicher.

\textbf{4.~``Rote'' Produkte als Ausnahme betrachten, nicht als Regel.} Es
gen\"{u}gt, ultra-verarbeitete Produkte zu reduzieren. Setzen Sie sich ein Limit
--- zum Beispiel maximal 1--2 ``rote'' Produkte pro Woche.

\section{\"{U}bersichtstabelle}

{\small
\begin{tabularx}{\textwidth}{|L|L|}
\hline
\rowcolor{verdescuro}\textcolor{white}{\textbf{Ergebnis in der App}}
  & \textcolor{white}{\textbf{Was in der Praxis tun}} \\
\hline
Wenige Zutaten, wenige/keine E-Nummern
  & T\"{a}gliche Basis: Naturjoghurt, Vollkornbrot, Tomatenpassata \\
\hline
Einige E-Nummern, moderater Zuckergehalt
  & Gelegentlich verwenden, nicht als Hauptnahrungsmittel \\
\hline
Viele E-Nummern, Zuckerzus\"{a}tze, komplexe Aromen
  & Einschr\"{a}nken oder meiden; nur als seltene Ausnahme \\
\hline
\end{tabularx}
}

\section{Warum das Lesen von Zusatzstoffen wirklich einen Unterschied macht}

Stark verarbeitete Lebensmittel zu erkennen ist keine ``Zusatzstoff-Phobie'' ---
es ist \textbf{Aufmerksamkeit f\"{u}r die Qualit\"{a}t dessen, was wir essen}.
Zahlreiche Studien legen nahe, dass die Reduzierung von Ultra-Verarbeitetem
langfristig bei der Gewichtskontrolle, der Herzgesundheit und der Stabilit\"{a}t
des Stoffwechsels hilft.

% ============================================================
\chapter{So lesen Sie ein N\"{a}hrwertetikett}

Zutaten, Kalorien, verstecktes Salz: lernen Sie, ein Etikett in unter einer
Minute zu entschl\"{u}sseln, und h\"{o}ren Sie auf, Produkte zu kaufen, die
gesund wirken, aber es nicht sind.

\section{Die Zutatenliste: worauf man zuerst schauen sollte}

Die Zutaten sind gesetzlich in \textbf{absteigender Reihenfolge nach Gewicht}
aufgef\"{u}hrt. Die erste Zutat ist am wichtigsten.

\begin{itemize}
  \item Steht Zucker, Glukose-Fruktose-Sirup oder Palmfett an erster Stelle,
    ist das Produkt um diese Substanz herum konzipiert.
  \item Die L\"{a}nge der Zutatenliste signalisiert den Verarbeitungsgrad:
    selbst gebackenes Brot enth\"{a}lt etwa 4 Zutaten, Industriebrot
    kann 15 oder mehr aufweisen.
  \item Komplizierte Namen wie ``modifizierte St\"{a}rke'',
    ``Carboxymethylcellulose'' oder ``hydriert'' deuten auf industrielle
    Zus\"{a}tze hin.
\end{itemize}

\boxinfo{Faustregel}{Wenn Sie nicht alle Zutaten in einem normalen
Lebensmittelgesch\"{a}ft finden w\"{u}rden, ist das Produkt wahrscheinlich
stark verarbeitet.}

\section{Die N\"{a}hrwerttabelle: worauf man achten sollte}

\begin{description}
  \item[Zucker] Unter 5\,g/100\,g = niedrig; 5--22,5\,g = mittel; \"{u}ber
    22,5\,g = hoch. F\"{u}r Getr\"{a}nke: unter 2,5\,g/100\,ml gilt als
    niedrig.
  \item[Salz] Unter 0,3\,g/100\,g = niedrig; \"{u}ber 1,5\,g/100\,g = hoch.
    Der auf dem Etikett angegebene Natriumwert muss mit 2,5 multipliziert
    werden, um den Salzgehalt zu erhalten.
  \item[Ges\"{a}ttigte Fetts\"{a}uren] \"{U}ber 5\,g/100\,g gilt als hoch.
\end{description}

\section{Portionen vs.\ 100\,g: der Trick, der verwirrt}

Hersteller pr\"{a}sentieren oft Portionswerte, die absichtlich niedrig
angesetzt sind. Nutzen Sie immer die Werte pro 100\,g f\"{u}r faire
Produktvergleiche.

\boxgelb{Beispiel}{Wenn die Packung 10\,g Zucker pro Portion (30\,g) angibt und
Sie 60\,g essen, nehmen Sie 20\,g Zucker zu sich --- fast die H\"{a}lfte der
Tagesempfehlung.}

\section{Verstecktes Salz: wo es sich wirklich verbirgt}

Salz erscheint auch in scheinbar s\"{u}{\ss}en Produkten wie Keksen,
M\"{u}sli und Brot.

\begin{itemize}
  \item \textbf{Brot und Backwaren:} zwei Scheiben Industriebrot enthalten
    0,5--0,8\,g Salz.
  \item \textbf{Fr\"{u}hst\"{u}cksflocken:} selbst ``gesunde'' und ``Vollkorn''-
    Varianten k\"{o}nnen 0,8--1,2\,g Salz pro 100\,g enthalten.
  \item \textbf{So{\ss}en und W\"{u}rzmittel:} zwischen 1 und 3\,g Salz
    pro 100\,g.
\end{itemize}

\section{N\"{a}hrwertbezogene Angaben: was sie wirklich bedeuten}

\textbf{``Ohne Zuckerzusatz'':} bedeutet nicht zuckerfrei; kann nat\"{u}rliche
Zucker enthalten. Nicht gleichbedeutend mit ``gesund''.

\textbf{``Light'' oder ``reduziert'':} bedeutet um 30\,\% reduziert; das
entfernte Fett wird oft durch Zucker oder Verdickungsmittel ersetzt.

\textbf{``Ballaststoffreich'' oder ``Proteinquelle'':} erfordert gesetzliche
Mindestwerte, sagt aber nichts \"{u}ber andere Inhaltsstoffe aus.

\section{Schnell-Checkliste f\"{u}r den Supermarkt}

\begin{enumerate}
  \item Erste Zutat ansehen: echtes Lebensmittel oder Zucker/verarbeitetes
    Fett?
  \item Zutaten z\"{a}hlen: weniger als 5 = ausgezeichnet; 5--10 = akzeptabel;
    \"{u}ber 10--12 = genaue Pr\"{u}fung erforderlich.
  \item E-Nummern pr\"{u}fen.
  \item Salz und Zucker pro 100\,g ansehen.
  \item Vorderseiten-Marketing ignorieren; der Zutatenliste vertrauen.
\end{enumerate}

\boxgruen{Abschlie{\ss}ender Rat}{Wenn Sie alle Zutaten erkennen und sie zu
Hause kochen k\"{o}nnten, sind Sie auf dem richtigen Weg.}

% ============================================================
\chapter{Zus\"{a}tze und Ultra-Verarbeitetes: Was ist der Unterschied?}

Ist ein Produkt mit vielen Zus\"{a}tzen zwangsl\"{a}ufig schlecht? Und ist
ultra-verarbeitetes immer voll von Zus\"{a}tzen? Die Antwort \"{u}berrascht.

\section{Was sind Lebensmittelzusatzstoffe}

Lebensmittelzusatzstoffe sind Stoffe, die gezielt hinzugef\"{u}gt werden, um
technologische Funktionen zu erf\"{u}llen: Haltbarkeit, F\"{a}rbung,
Verdickung, Emulgierung, Stabilisierung oder Geschmacksverbesserung. In Europa
erh\"{a}lt jeder zugelassene Zusatzstoff einen E-Code und muss eine
Sicherheitsbewertung der EFSA bestehen.

\begin{itemize}
  \item Nicht alle Zusatzstoffe sind synthetisch: Zitronens\"{a}ure, Curcumin,
    Sojalecithin und Ascorbins\"{a}ure sind nat\"{u}rlichen Ursprungs.
  \item Die meisten zugelassenen E-Codes gelten in erlaubten Dosierungen
    als sicher.
  \item Das Problem liegt in der Menge und Kombination verschiedener
    Zusatzstoffe.
\end{itemize}

\section{Was sind ultra-verarbeitete Lebensmittel}

Der Begriff stammt aus der NOVA-Klassifizierung der Universit\"{a}t S\~{a}o
Paulo. Ein Lebensmittel gilt als ultra-verarbeitet (NOVA~4), wenn es Zutaten
enth\"{a}lt, die in einer Haushaltsk\"{u}che nicht vorkommen: isolierte Proteine,
modifizierte St\"{a}rken, Proteinhydrolysate, Glukose-Fruktose-Sirupe,
``naturidentische'' Aromen, Emulgatoren, Farbstoffe und Stabilisatoren.

\textbf{Ultra-verarbeitet:} Fertigsnacks, S\"{u}{\ss}igkeiten, gezuckerte
Fr\"{u}hst\"{u}cksflocken, Softdrinks, rekonstituierte W\"{u}rstchen,
tiefgefrorene Fertiggerichte, aromatisierter Joghurt, industrielles Toastbrot.

\textbf{KEIN ultra-verarbeitetes Produkt:} gereifter K\"{a}se, qualit\"{a}tsvoller
Aufschnitt, handwerkliches Brot, traditionelle Pasta, Tomatenkonserven ohne
Zus\"{a}tze, Essig und natives Oliven\"{o}l.

\section{Wann sie \"{u}bereinstimmen --- und wann nicht}

\begin{description}
  \item[Zus\"{a}tze, aber KEIN Ultra-Verarbeitetes:] K\"{a}se mit
    Konservierungsmittel E252, Wein mit Sulfiten E220, \"{O}l mit
    Vitamin~E E306.
  \item[Ultra-verarbeitet mit WENIGEN E-Zus\"{a}tzen:] Toastbrot mit
    modifizierter St\"{a}rke und Aromen, Proteinriegel, Surimi.
  \item[Ultra-verarbeitet UND viele Zus\"{a}tze:] Industriegeb\"{a}ck,
    Industrieso{\ss}en, Energy Drinks.
\end{description}

\boxinfo{Fazit}{Ein Produkt mit einem nat\"{u}rlichen~E ist kein Problem. Ein
Produkt mit 15 industriellen Zutaten ohne E-Codes ist trotzdem
ultra-verarbeitet.}

\section{Das doppelte Problem: warum beides z\"{a}hlt}

\textbf{Risiken durch spezifische Zusatzstoffe:} Azofarbstoffe werden mit
Hyperaktivit\"{a}t bei Kindern verbunden; Nitrate in verarbeitetem Fleisch
mit krebserregenden Nitrosaminen; intensive S\"{u}{\ss}ungsmittel mit
m\"{o}glichen Ver\"{a}nderungen des Darmmikrobioms.

\textbf{Risiken durch Ultra-Verarbeitung:} Studien wie NutriNet-Sant\'{e} und
UK Biobank verbinden hohen Konsum von NOVA~4 mit erh\"{o}htem Risiko f\"{u}r
Adipositas, Typ-2-Diabetes, Herz-Kreislauf-Erkrankungen und bestimmten
Krebsarten.

\textbf{Der Kombinationseffekt:} Im echten Leben werden 10--20 verschiedene
Zus\"{a}tze kombiniert konsumiert. Die synergistischen Effekte sind noch
wenig erforscht.

\section{Wie man sie in der Praxis erkennt}

\begin{enumerate}
  \item Pr\"{u}fen auf ``Laborzutaten'': modifizierte St\"{a}rken, isolierte
    Proteine, Hydrolysate, Glukose-Fruktose-Sirupe.
  \item E-Codes z\"{a}hlen und kritische Kategorien identifizieren.
  \item Den Mahlzeitenkontext ber\"{u}cksichtigen.
  \item Die App als Unterst\"{u}tzung nutzen, nicht als absoluten Richter.
\end{enumerate}

% ============================================================
\chapter{NOVA-Klassifikation: Was ist sie und wie verwendet man sie?}

NOVA misst keine N\"{a}hrstoffe, sondern den \emph{Grad der industriellen
Verarbeitung}. Die 4 Gruppen zu verstehen ver\"{a}ndert die Art, wie Sie
einkaufen.

\section{Was ist die NOVA-Klassifizierung}

NOVA ist ein Klassifizierungssystem f\"{u}r Lebensmittel, das von der
Forschungsgruppe von Professor Carlos Monteiro an der Universit\"{a}t S\~{a}o
Paulo entwickelt wurde und heute von der WHO, der FAO und zahlreichen
epidemiologischen Studien weltweit verwendet wird.

Im Gegensatz zu n\"{a}hrstoffbasierten Systemen wie dem Nutri-Score
klassifiziert NOVA nach dem \emph{Grad der industriellen Verarbeitung}:
nicht die Kalorienzahl, sondern den Verarbeitungsgrad ist entscheidend.

Industrieprozesse k\"{o}nnen Faserstruktur, N\"{a}hrstoffverf\"{u}gbarkeit,
Darmmikrobiom und S\"{a}ttigungsreaktion ver\"{a}ndern --- unabh\"{a}ngig vom
finalen N\"{a}hrstoffgehalt. NOVA unterteilt alle Lebensmittel in vier
verschiedene kategorische Gruppen.

\section{Gruppe 1 --- Unverarbeitete oder minimal verarbeitete Lebensmittel}

Lebensmittel ohne wesentliche industrielle Verarbeitung oder nur minimale
physikalische Prozesse: Trocknen, Mahlen, K\"{u}hlen, Pasteurisieren.

\textbf{Beispiele:} frisches oder tiefgefrorenes Obst und Gem\"{u}se, frisches
oder tiefgefrorenes Fleisch und Fisch ohne Zus\"{a}tze, Eier, getrocknete
H\"{u}lsenfr\"{u}chte, Vollkorngetreide (Reis, Hafer, Dinkel), Vollkornmehl,
frische pasteurisierte Milch, Naturjoghurt ohne Zus\"{a}tze, Tee, Kaffee.

\boxgruen{Praktisches Ziel}{Lebensmittel der Gruppe~1 sollten den gr\"{o}{\ss}ten
Teil jeder Mahlzeit ausmachen.}

\section{Gruppe 2 --- Verarbeitete K\"{u}chenzutaten}

Substanzen aus Gruppe-1-Lebensmitteln oder der Natur, die in der K\"{u}che zum
Kochen und W\"{u}rzen verwendet werden. Sie werden nicht isoliert konsumiert.

\textbf{Beispiele:} Oliven\"{o}l und pflanzliche \"{O}le, Butter, Essig, Salz,
Zucker, Honig, Weizenmehl (Typ~405), reine Maisst\"{a}rke, Ahornsirup,
Milchpulver und traditionelle H\"{a}rtk\"{a}sesorten.

M\"{a}{\ss}ig verwenden: \"{u}berm\"{a}{\ss}iger Zucker, Salz und ges\"{a}ttigte
Fette bleiben ung\"{u}nstig.

\section{Gruppe 3 --- Verarbeitete Lebensmittel}

Produkte durch Hinzuf\"{u}gen von Gruppe-2-Zutaten zu Gruppe-1-Lebensmitteln,
durch Verfahren wie Einsalzen, R\"{a}uchern, Fermentieren oder Reifen.
Sie enthalten oft nur 2--3 Zutaten.

\textbf{Beispiele:} konserviertes Gem\"{u}se und H\"{u}lsenfr\"{u}chte (mit
Wasser und Salz), kandierte und eingelegte Fr\"{u}chte, ger\"{o}stete und
gesalzene N\"{u}sse, Sardellen und Sardinen in \"{O}l, qualit\"{a}tsvoller
roher Schinken, traditionelle K\"{a}sesorten (Parmesan, Pecorino), handwerkliches
Brot, Wein, handwerkliches Bier.

Gruppe~3 muss nicht vermieden werden. Viele traditionelle Lebensmittel
geh\"{o}ren hierher.

\section{Gruppe 4 --- Ultra-Verarbeitetes: die Gruppe, die man einschr\"{a}nken
sollte}

Ultra-verarbeitete Lebensmittel sind industrielle Zubereitungen ohne ganze
Lebensmittel oder nur minimal. Sie bestehen aus extrahierten, gereinigten
oder chemisch ver\"{a}nderten Zutaten sowie gro{\ss}en Mengen technologischer
Zus\"{a}tze.

Das Schl\"{u}sselsignal: Sie enthalten \textbf{Zutaten, die in einer
Haushaltsk\"{u}che nicht vorkommen}.

Typische Zutaten: hydrolysierte Proteine, modifizierte St\"{a}rken,
Glukose-Fruktose-Sirup, hydrierte \"{O}le, ``naturidentische'' Aromen,
k\"{u}nstliche Farbstoffe, intensive S\"{u}{\ss}ungsmittel (Aspartam, Saccharin,
Sucralose), Emulgatoren (E471, E472).

\boxrot{Wichtige Daten}{In Deutschland stammen etwa 15--20\,\% der t\"{a}glichen
Kalorien aus ultra-verarbeiteten Lebensmitteln. Ein Konsum von mehr als 20\,\%
der t\"{a}glichen Kalorien aus NOVA~4 ist mit Gesundheitsrisiken verbunden.}

\section{NOVA beim Einkaufen in der Praxis}

\textbf{1.~``K\"{o}nnte ich das zu Hause mit normalen Zutaten kochen?''}
Eine bejahende Antwort deutet auf Gruppe~1, 2 oder~3. Eine Verneinung bedeutet
fast sicher Gruppe~4.

\textbf{2.~Nach den ``Signalw\"{o}rtern'' der Gruppe~4 suchen:}
``Modifizierte St\"{a}rke'', ``isoliertes Sojaprotein'', ``Maltodextrine'',
``Sirup aus\ldots'', ``komplexe Aromen'', ``hydriertes Palmfett'',
``Natriumcaseinat''.

\textbf{3.~Praktische Regel:} mehr als die H\"{a}lfte des Tellers aus Gruppe~1.

\section{Vergleichstabelle NOVA}

{\small
\begin{tabularx}{\textwidth}{|l|L|L|l|}
\hline
\rowcolor{verdescuro}
  \textcolor{white}{\textbf{Gruppe}}
  & \textcolor{white}{\textbf{Definition}}
  & \textcolor{white}{\textbf{Beispiele}}
  & \textcolor{white}{\textbf{Ern\"{a}hrung}} \\
\hline
NOVA~1 & Nicht/min. verarbeitet & Obst, Gem\"{u}se, Eier, frisches Fleisch & Basis \\
\hline
NOVA~2 & K\"{u}chenzutaten & \"{O}l, Salz, Zucker, Mehl, Butter & M\"{a}{\ss}ig \\
\hline
NOVA~3 & Trad. verarbeitet & Konserven, K\"{a}se, Rohschinken, Brot & Akzeptabel \\
\hline
NOVA~4 & Ultra-verarbeitet & Geb\"{a}ck, W\"{u}rstchen, gezuck. Flocken & Einschr\"{a}nken \\
\hline
\end{tabularx}
}

% ============================================================
\chapter{Zutaten, auf die man bei Etiketten achten sollte}

Sie m\"{u}ssen kein Chemiker sein, um Etiketten richtig zu lesen. Mit einigen
Schl\"{u}sselbegriffen k\"{o}nnen Sie problematische Produkte vermeiden und
schnell bessere Entscheidungen treffen.

\section{Zucker: die 10 versteckten Namen auf Etiketten}

Zucker wird unter vielen verschiedenen Namen aufgef\"{u}hrt, um seine
Bedeutung in der Zutatenliste zu verschleiern. Die 10 h\"{a}ufigsten Synonyme:

\begin{itemize}
  \item Glukose-Fruktose-Sirup
  \item Maltodextrine
  \item Dextrose
  \item Fruktose
  \item Konzentrierter Rohrzuckersaft
  \item Maissirup / Reissirup
  \item Getrockneter Honig
  \item Kokoszucker
  \item Agavennektar
\end{itemize}

\boxgelb{Die Falle ``ohne Zuckerzusatz''}{Ein Saft mit dieser Kennzeichnung
kann trotzdem 10--12\,g Zucker pro 100\,ml enthalten --- mehr als Cola. Die
WHO empfiehlt, freie Zucker unter 10\,\% der t\"{a}glichen Kalorien zu halten
(etwa 50\,g/Tag f\"{u}r einen Erwachsenen mit 2000\,kcal), vorzugsweise
unter 5\,\% (25\,g).}

\section{Fette, die man meiden sollte: hydrierte Fette und Transfetts\"{a}uren}

\textbf{Hydrierte und teilweise hydrierte Fette:} auf Etiketten nach Begriffen
wie ``teilweise hydriertes Pflanzen\"{o}l'', ``hydriertes Pflanzenfett'' oder
``Shortening'' suchen. In Europa sind diese seit 2021 \"{u}ber 2\,g/100\,g
Fett verboten.

\textbf{Palmfett:} ges\"{a}ttigt, aber nicht hydriert. Das Risiko liegt in
m\"{o}glicherweise krebserregenden Verbindungen (3-MCPD, Glycidol), die bei
der Raffination bei hohen Temperaturen entstehen.

\boxgruen{Unbedenkliche Fette}{Natives Oliven\"{o}l extra, Lein\"{o}l, N\"{u}sse
und \"{O}lsaaten, Avocado, Fette aus fettem Fisch.}

\section{Umstrittene Zus\"{a}tze: die E-Codes, die man kennen sollte}

\textbf{Azofarbstoffe --- besondere Vorsicht bei Kindern:}
E102 (Tartrazin), E104 (Chinolgelb), E110 (Gelborange~S), E122 (Azorubin),
E124 (Ponceau 4R), E129 (Allurarot~AC). Eine EFSA-Studie von 2007 zeigte einen
Zusammenhang mit Hyperaktivit\"{a}t bei Kindern. In Europa ist ein Warnhinweis
auf dem Etikett Pflicht.

\textbf{Nitrate und Nitrite (E249--E252):} Konservierungsstoffe in Aufschnitt,
W\"{u}rstchen, Bacon. In Gegenwart tierischer Proteine entstehen Nitrosamine,
die als wahrscheinlich krebserregend (IARC Gruppe~2A) eingestuft sind.

\textbf{Intensive S\"{u}{\ss}ungsmittel (E950, E951, E952, E954, E955):}
in zugelassenen Dosen f\"{u}r Erwachsene sicher, aber die WHO (2023) \"{a}u{\ss}erte
Vorsicht beim chronischen Gebrauch zur Gewichtskontrolle.

\section{Verstecktes Natrium: mehr als nur Kochsalz}

Etwa 75\,\% des Natriums in der westlichen Ern\"{a}hrung stammt nicht aus dem
Kochen, sondern ist bereits in verpackten Produkten enthalten.

Wichtigste Quellen: Brot, K\"{a}se, Aufschnitt, Fertiggerichte, So{\ss}en im
Glas, W\"{u}rfelbriihe, Sojasosse, Ketchup, Mayonnaise, salzige Snacks,
industrielle Fr\"{u}hst\"{u}cksflocken.

\boxinfo{Umrechnung Natrium in Salz}{Den Natriumwert mit 2,5 multiplizieren.
Beispiel: 0,6\,g Natrium = 1,5\,g Salz. WHO-Grenze: 5\,g Salz pro Tag
(etwa 2\,g Natrium).}

\section{Die 5er-Regel: der schnellste Filter}

\begin{enumerate}
  \item \textbf{Mehr als 5 Zutaten?} Dann genauer lesen.
  \item \textbf{Zucker \"{u}ber 5\,g/100\,g?} Pr\"{u}fen, ob er an erster
    oder zweiter Stelle steht.
  \item \textbf{Salz \"{u}ber 0,5\,g/100\,g in einem s\"{u}{\ss}en Produkt?}
    Signal f\"{u}r Ultra-Verarbeitung.
  \item \textbf{Mehr als 5 E-Zus\"{a}tze?} Das Produkt ist wahrscheinlich
    ultra-verarbeitet.
  \item \textbf{5 oder mehr Zutaten nicht erkannt?} Zur\"{u}ck ins Regal damit.
\end{enumerate}

\boxinfo{Beachte}{Die 5er-Regel ist eine Orientierung, kein Urteil. Verwende
sie als Ausgangspunkt.}

% ============================================================
\chapter{Empfohlene Lebensmittel: Die gesundesten Optionen im Supermarkt}

Es geht nicht darum, perfekt zu sein oder nur Salat zu essen. Es geht darum
zu wissen, welche Produkte bei gleicher Bequemlichkeit wirklich besser sind
als andere.

\section{Das Grundprinzip: ersetzen, nicht eliminieren}

\begin{description}
  \item[Gleicher Aufwand, weniger Verarbeitung.] Naturjoghurt statt
    aromatisiertem Joghurt, Vollkornreis statt vorgekochtem Reis,
    Konserven-H\"{u}lsenfr\"{u}chte statt Fertiggerichte.
  \item[Das 80/20-Prinzip funktioniert.] Wenn 80\,\% der Nahrung aus weniger
    verarbeiteten Gruppen stammt, beeintr\"{a}chtigen st\"{a}rker verarbeitete
    Produkte (20\,\%) die Gesamtgesundheit nicht wesentlich.
  \item[``Empfohlene'' Lebensmittel sind nicht alle Bio oder teuer.]
    Konserven-Kichererbsen, Eier, Vollkornmehl, Naturjoghurt, Haferflocken,
    Linsen kosten oft weniger als verpackte Alternativen.
\end{description}

\section{Getreide und Kohlenhydrate: die besten Entscheidungen}

\textbf{Empfohlene Getreidesorten:} Haferflocken (ohne Zus\"{a}tze),
Vollkornreis, Dinkelvollkorn, Graupen, Buchweizen, Vollkornnudeln aus reinem
Vollkorngriei{\ss}, Brot aus reinem Vollkornmehl (maximal 4--5 Zutaten).

\textbf{Einzuschr\"{a}nken:} gezuckerte Fr\"{u}hst\"{u}cksflocken (auch
``Vollkorn''-Varianten mit Sirup), Toastbrot mit modifizierten St\"{a}rken,
Cracker mit Palmfett, vorgekochter Reis und Nudeln im Beutel mit Zus\"{a}tzen.

\boxgelb{Schnelltipp}{Wenn der Zuckergehalt \"{u}ber 10\,g/100\,g liegt, ist
das ein Dessertprodukt --- kein gesundes Fr\"{u}hst\"{u}ck.}

\section{Proteine: die besten Quellen}

\textbf{Fisch und Meeresfr\"{u}chte:} frischer oder tiefgefrorener Fisch
ohne Panade, Sardinen und Makrelen in Oliven\"{o}l, Thunfisch im eigenen
Saft, Sardellen in Salz. Fettreicher Fisch z\"{a}hlt zu den g\"{u}nstigsten
und n\"{a}hrstoffreichsten Proteinquellen.

\textbf{Eier und Milchprodukte:} ganze Eier, griechischer Naturjoghurt,
Ricotta und H\"{u}tttenk\"{a}se. Zu meiden sind industrielle Eiersatze,
``Protein''-Joghurts mit Verdickungsmitteln, Schmelzk\"{a}se mit Phosphaten.

\textbf{H\"{u}lsenfr\"{u}chte:} Linsen, Kichererbsen, Bohnen, Erbsen und
Edamame. Dosenvarianten im eigenen Saft sind g\"{u}nstig, ballaststoffreich
und haben ein ausgezeichnetes Aminos\"{a}ureprofil.

\section{Gute Fette: die einfachsten zum Einbauen}

Qualit\"{a}tsfette wie natives Oliven\"{o}l extra, N\"{u}sse, Avocado und
fetter Fisch geh\"{o}ren zu den herzsch\"{u}tzendsten Komponenten.

\textbf{Empfohlen:} natives Oliven\"{o}l extra, rohes Lein\"{o}l, Waln\"{u}sse,
Mandeln, Haseln\"{u}sse, Leinsamen, Chiasamen, Avocado, fetter Fisch.

\textbf{In Ma{\ss}en:} Butter, Kokosn\"{o}ss\"{o}l, Vollmilch.

\textbf{Fette, die man meiden sollte:} Margarine mit teilweise hydrierten
Fetten, pflanzliches Shortening, raffiniertes Palmfett, Pflanzencreme mit
Emulgatoren.

\section{Clevere Snacks}

\begin{enumerate}
  \item \textbf{N\"{u}sse und Saaten natur:} Waln\"{u}sse, Mandeln, Haseln\"{u}sse,
    Cashews ohne Salzzusatz (30\,g liefern Proteine, gute Fette und Mineralstoffe).
  \item \textbf{Frisches ganzes Obst:} enth\"{a}lt Ballaststoffe, die die
    Zuckeraufnahme verlangsamen.
  \item \textbf{Griechischer Naturjoghurt:} mit hohem Proteingehalt, niedrigem
    glyk\"{a}mischen Index und ohne Zus\"{a}tze.
  \item \textbf{H\"{a}rtk\"{a}se + Vollkorncracker:} S\"{a}ttigung durch
    einfache Zutaten.
\end{enumerate}

\section{Muster-Einkaufsliste}

{\small
\begin{tabularx}{\textwidth}{|l|L|}
\hline
\rowcolor{verdescuro}
  \textcolor{white}{\textbf{Kategorie}}
  & \textcolor{white}{\textbf{Empfohlene Produkte}} \\
\hline
Getreide & Haferflocken, Vollkornreis, Dinkel, Vollkornnudeln, Vollkornbrot \\
\hline
H\"{u}lsenfr\"{u}chte & Linsen, Kichererbsen, Bohnen, Erbsen \\
\hline
Tierische Proteine & Eier, Sardinen, Thunfisch, griechischer Naturjoghurt \\
\hline
Gem\"{u}se & Saisonal, frisch oder tiefgek\"{u}hlt ohne Zus\"{a}tze \\
\hline
Fette & Natives Oliven\"{o}l extra, Waln\"{u}sse, Mandeln, Leinsamen \\
\hline
Obst & Saisonales ganzes Obst, tiefgefrorene Beeren (ohne Zucker) \\
\hline
W\"{u}rzmittel & Apfelessig, wenig verarbeitete Sojasosse, Gew\"{u}rze \\
\hline
\end{tabularx}
}

\boxgruen{Der Warenkorb-Test}{Beim n\"{a}chsten Einkauf z\"{a}hlen, wie viele
Produkte weniger als 5 Zutaten haben. Wenn es mehr als die H\"{a}lfte sind,
sind Sie auf dem richtigen Weg.}

% ============================================================
\chapter{Warum ultra-verarbeitete Lebensmittel sch\"{a}dlich sind}

Es geht nicht nur um Kalorien oder Fett: Ultra-verarbeitete Lebensmittel
st\"{o}ren das S\"{a}ttigungsgef\"{u}hl, ver\"{a}ndern das Mikrobiom und
l\"{o}sen chronische Entz\"{u}ndungen aus.

\section{Fettleibigkeit und Stoffwechselst\"{o}rungen}

Ultra-verarbeitete Lebensmittel sind so entwickelt, dass sie \emph{\"{u}beraus
attraktiv} schmecken und aktiv in die nat\"{u}rlichen S\"{a}ttigungssignale
des K\"{o}rpers eingreifen.

\boxinfo{Kevin Hall Studie (NIH, 2019)}{Teilnehmer mit ultra-verarbeiteter
Ern\"{a}hrung nahmen \emph{im Durchschnitt 500 Kalorien mehr pro Tag} zu sich
als jene mit minimal verarbeiteten Lebensmitteln --- trotz gleicher
Verf\"{u}gbarkeit. In zwei Wochen nahm die UPF-Gruppe fast ein Kilogramm zu.}

UPF enthalten im Durchschnitt 2,15\,kcal pro Gramm --- etwa doppelt so viel
wie frische Lebensmittel. Sie reduzieren das Peptid~YY, das Darmhormon, das
dem Gehirn S\"{a}ttheit signalisiert.

Statistik: \textbf{+41\,\% Risiko f\"{u}r abdominale Fettleibigkeit} bei
Personen mit hohem UPF-Konsum.

\section{Entz\"{u}ndungen und Dysfunktionen des Immunsystems}

H\"{a}ufiger Konsum ultra-verarbeiteter Lebensmittel korreliert mit
verst\"{a}rkten entz\"{u}ndlichen Immunreaktionen. Bestimmte in UPF h\"{a}ufig
vorkommende Substanzen --- Emulgatoren, Verdickungsmittel und Titandioxid
(E171) --- k\"{o}nnen das Darmmikrobiom ver\"{a}ndern und die
Durchl\"{a}ssigkeit der Darmschleimhaut erh\"{o}hen (\textit{leaky gut}).

\textbf{Mit hohem UPF-Konsum assoziierte Autoimmunerkrankungen:}
Z\"{o}liakie, Hashimoto-Thyreoiditis, Multiple Sklerose, systemischer Lupus
erythematodes, Typ-1-Diabetes.

\section{Darmgesundheit und Mikrobiom}

Der Ultraverarbeitungsprozess entzieht Lebensmitteln Ballaststoffe und
bioaktive pflanzliche Verbindungen, die zur Ern\"{a}hrung n\"{u}tzlicher
Darmbakterien notwendig sind.

\textbf{Dysbiose:} ein verarmtes Mikrobiom ist mit Insulinresistenz, chronischer
Entz\"{u}ndung und erh\"{o}hter Anf\"{a}lligkeit f\"{u}r Infektionen verbunden.

\textbf{H\"{a}ufige Erkrankungen:} Morbus Crohn, Colitis ulcerosa,
Reizdarmsyndrom, Magengeschw\"{u}re, pr\"{a}kanzer\"{o}se
Dickdarmpolypen.

\section{Ein systemisches Risiko: mehrere Mechanismen gleichzeitig}

{\small
\begin{tabularx}{\textwidth}{|l|L|L|}
\hline
\rowcolor{verdescuro}
  \textcolor{white}{\textbf{Mechanismus}}
  & \textcolor{white}{\textbf{Wirkung}}
  & \textcolor{white}{\textbf{Langzeitfolge}} \\
\hline
Hyperattraktivit\"{a}t & \"{U}berm\"{a}{\ss}ige Kalorienaufnahme & Fettleibigkeit, Typ-2-Diabetes \\
\hline
Ballaststoffmangel & Intestinale Dysbiose & Chronische Entz\"{u}ndung \\
\hline
Zus\"{a}tze (Emulgatoren) & Darmdurchl\"{a}ssigkeit & Autoimmun-, Herz-Kreislauf-Erkrankungen \\
\hline
Zucker und S\"{u}{\ss}ungsmittel & Insulinresistenz & Typ-2-Diabetes, metabolisches Syndrom \\
\hline
Verdr\"{a}ngung frischer Lebensmittel & Mikron\"{a}hrstoffmangel & Oxidativer Stress \\
\hline
\end{tabularx}
}

Gro{\ss}angelegte Studien --- NutriNet-Sant\'{e} (\"{u}ber 100.000 Teilnehmer)
und UK Biobank (\"{u}ber 500.000 Teilnehmer) --- verbinden hohen UPF-Konsum
mit erh\"{o}hten Risiken f\"{u}r Krebs, Herz-Kreislauf-Erkrankungen,
Depressionen und Gesamtmortalit\"{a}t.

\section{Was man in der Praxis tun kann}

Die praktische Botschaft lautet nicht ``eliminiere alles''. Sie lautet:
\emph{Reduziere den Anteil von UPF in der t\"{a}glichen Ern\"{a}hrung}.

\begin{enumerate}
  \item \textbf{Identifiziere deine gew\"{o}hnlichen UPF:} gezuckerte
    Fr\"{u}hst\"{u}ckszerealien, abgepackte Snacks, Toastbrot, Fertiggerichte,
    aromatisierte Getr\"{a}nke.
  \item \textbf{Ersetze eines nach dem anderen:} Haferflocken statt
    s\"{u}{\ss}er Cerealien; Naturjoghurt statt aromatisiertem mit
    Verdickungsmitteln; Handwerksbrot statt Toastbrot.
  \item \textbf{Erh\"{o}he die Ballaststoffzufuhr --- Priorit\"{a}t:}
    H\"{u}lsenfr\"{u}chte, Gem\"{u}se, Vollkorngetreide, ganzes Obst.
  \item \textbf{Nutze die App, um versteckte UPF zu erkennen:}
    E-Codes Reader zeigt NOVA-Klassifikation und Zus\"{a}tze sofort.
  \item \textbf{Strebe keine Perfektion an:} Selbst eine Verbesserung um
    20--30\,\% erzeugt mittelfristig messbare Auswirkungen.
\end{enumerate}

\boxgruen{Gute Nachricht}{Selbst eine Verbesserung der Ern\"{a}hrungsqualit\"{a}t
um 20--30\,\% erzeugt mittelfristig messbare gesundheitliche Auswirkungen.}

% ============================================================
\chapter{Lebensmittelsicherheit f\"{u}r Kinder und \"{a}ltere Menschen}

Kinder und \"{a}ltere Menschen reagieren auf Zusatzstoffe und Zutaten minderer
Qualit\"{a}t anders als gesunde Erwachsene.

\section{Warum Kinder und \"{a}ltere Menschen anf\"{a}lliger sind}

\textbf{Kinder: geringes K\"{o}rpergewicht, unreife Systeme.}
Ein 15\,kg schweres Kind erh\"{a}lt bei gleichem Produktkonsum eine relative
Dosis, die 4--5 Mal h\"{o}her ist als die eines 70\,kg schweren Erwachsenen.
Leber und Nieren metabolisieren bestimmte Substanzen langsamer.

\textbf{\"{A}ltere Menschen: eingeschr\"{a}nkte Nieren- und Leberfunktion.}
Mit dem Alter sinkt die Ausscheidungsf\"{a}higkeit. Mehrfachmedikation kann
mit Zus\"{a}tzen wechselwirken, und \"{u}berschussiges Natrium belastet
Menschen mit Bluthochdruck st\"{a}rker.

\textbf{Kumulative Effekte:}
Einzelne Portionen sind oft unbedenklich, doch regelm\"{a}{\ss}iger Konsum
\"{u}ber den Tag verteilt f\"{u}hrt zu relevanten kumulativen Belastungen.

\section{Kritische Zus\"{a}tze f\"{u}r Kinder}

\textbf{Azofarbstoffe --- bei Kindern zu meiden:}
E102 (Tartrazin), E104 (Chinolgelb), E110 (Gelborange~S~FCF), E122 (Azorubin),
E124 (Ponceau 4R), E129 (Allurarot~AC).
In Europa ist auf dem Etikett folgender Hinweis Pflicht: ``kann die Aktivit\"{a}t
und Aufmerksamkeit von Kindern beeintr\"{a}chtigen''.

\textbf{Intensive S\"{u}{\ss}ungsmittel --- nicht f\"{u}r Kleinkinder
zugelassen:}
E951, E950, E955, E954. In Lebensmitteln f\"{u}r S\"{a}uglinge und Kleinkinder
nicht zugelassen (EU-Verordnung 1333/2008). Pr\"{a}sent in zuckerfreien
Joghurts, \textit{Light}-Getr\"{a}nken, Hustensirups, zuckerfreien Bonbons.

\boxrot{Koffein}{Die EFSA empfiehlt maximal 3\,mg/kg/Tag. Eine 250-ml-Dose
Energy Drink kann bis zu 80\,mg Koffein enthalten --- fast das Doppelte der
Grenze f\"{u}r ein 15\,kg schweres Kind.}

\section{Was \"{a}ltere Menschen einschr\"{a}nken sollten}

\textbf{Natrium: das Hauptrisiko.} Bluthochdruck und Niereninsuffizienz sind
h\"{a}ufig bei \"{a}lteren Menschen. Ziel: weniger als 5\,g Salz t\"{a}glich.
Hauptvektoren: Fertiggerichte, Aufschnitt, H\"{a}rk\"{a}se, Instant-Br\"{u}hen.

\textbf{Phosphate (E338--E452):} k\"{o}nnen die Kalziumaufnahme st\"{o}ren ---
ein Problem f\"{u}r Menschen mit Osteoporose-Risiko. Kommen in verarbeiteten
Fleischprodukten und Schmelzk\"{a}se vor.

\textbf{Wechselwirkungen mit Medikamenten:} Vitamin~K (gr\"{u}nes Blattgem\"{u}se)
mit Cumarinen; Grapefruit mit Statinen; Cranberrysaft mit bestimmten
Antikoagulanzien.

\section{Produkte ``f\"{u}r Kinder'', bei denen man genauer hinschauen sollte}

\textbf{Fr\"{u}hst\"{u}cksflocken ``f\"{u}r Kinder'':} enthalten oft 25--35\,g
Zucker pro 100\,g, k\"{u}nstliche Farbstoffe und zugesetzte Vitamine.

\textbf{Getr\"{a}nke ``mit Fruchtsaft'':} k\"{o}nnen nur 10--30\,\% echten
Saft enthalten. Immer den Fruchtanteil \"{u}berpr\"{u}fen.

\textbf{Joghurt ``f\"{u}r Kinder'' mit Frucht:} enth\"{a}lt oft mehr Zucker
als Erwachsenenjoghurts, mehr Verdickungsmittel und k\"{u}nstliche Aromen.
Griechischer Naturjoghurt mit frischem Obst ist besser.

\section{Sichere Einkaufsliste}

\begin{enumerate}
  \item \textbf{F\"{u}r Kinder:} Produkte mit wenigen erkennbaren Zutaten
    w\"{a}hlen. Meiden: mehr als 2--3 Farbstoffe oder k\"{u}nstliche
    S\"{u}{\ss}ungsmittel.
  \item \textbf{Immer den Zuckergehalt in ``Kinder''-Produkten pr\"{u}fen:}
    Richtwert: weniger als 10\,g/100\,g.
  \item \textbf{F\"{u}r \"{a}ltere Menschen:} Priorit\"{a}t auf niedrigen
    Natriumgehalt; weniger als 0,3\,g Salz/100\,g.
  \item \textbf{Vor dem Kauf mit der App scannen:}
    E-Codes Reader zeigt sofort kritische Zus\"{a}tze und Natriumgehalt.
  \item \textbf{F\"{u}r Hauptmahlzeiten immer Gruppe~1--2 (NOVA) bevorzugen:}
    Gem\"{u}se, H\"{u}lsenfr\"{u}chte, Eier, Fisch, Naturjoghurt, ganzes Obst.
\end{enumerate}

\section{Vergleichstabelle}

{\small
\begin{tabularx}{\textwidth}{|L|l|l|}
\hline
\rowcolor{verdescuro}
  \textcolor{white}{\textbf{Zusatzstoff/Zutat}}
  & \textcolor{white}{\textbf{Kinder}}
  & \textcolor{white}{\textbf{\"{A}ltere Menschen}} \\
\hline
Azofarbstoffe (E102, E110, E122\ldots) & Meiden & Begrenzen \\
\hline
Intensive S\"{u}{\ss}ungsmittel & Nicht f\"{u}r Kleinkinder & In Ma{\ss}en \\
\hline
Hoher Natriumgehalt (>1,5\,g Salz/100\,g) & Begrenzen & Meiden \\
\hline
Phosphate (E338--E452) & M\"{a}{\ss}ig & Begrenzen \\
\hline
Koffein & Meiden <12 Jahre & Achtung \\
\hline
Nitrate/Nitrite (E249--E252) & Begrenzen & Begrenzen \\
\hline
\end{tabularx}
}

% ============================================================
\backmatter

\chapter*{Impressum und App-Download}
\addcontentsline{toc}{chapter}{Impressum}

\begin{center}
\vspace{1cm}
{\large\bfseries\color{verdescuro} Bewusst Essen}\\[0.4cm]
{\itshape Praktischer Leitfaden zu Lebensmittelzusatzstoffen}\\[0.3cm]
{\small Mit E-Codes Reader f\"{u}r Android}

\vspace{1.5cm}
\textbf{Version 1.0} \quad \textbullet \quad 2026\\
\textbf{Lizenz:} Creative Commons BY-NC 4.0\\[0.3cm]
\textbf{Website:} \href{https://faustobe.it}{faustobe.it}\\
\textbf{App:} \href{https://play.google.com/store/apps/details?id=it.faustobe.ecodes}{E-Codes Reader bei Google Play}

\vspace{1.5cm}
{\small QR-Code scannen, um die App herunterzuladen:}\\[0.5cm]

\qrcode[height=4cm]{https://play.google.com/store/apps/details?id=it.faustobe.ecodes}

\vspace{1cm}
{\footnotesize Dieses E-Book wurde mit \LaTeX{} von faustobe erstellt.}
\end{center}

\end{document}
